\chapter{Superficies y sus curvas de intersección}

En la sección \secref{regionesparametrizadas} se presentaron algunas generalidades
sobre parametrizaciones básicas.
\section{Curvas parametrizadas} \index{Curva!parametrizada}
De acuerdo a \cite{thomas_calculo_2015}, consideremos que $f$ y $g$ son funciones continuas en $[a,b]$.
Si $x=f(t) $ y $y=g(t)$, el conjunto de puntos $(x,y)$ que satisfacen estas ecuaciones se denomina curva paramétrica
(o parametrizada) en $\mathbb{R}^2$ y las ecuaciones que la definen se denominan las ecuaciones paramétricas
de la curva. Si además $f'(t)$ y $g'(t)$ no son simultaneamente cero, la curva se denomina suave.
\index{Curva!suave}
\section{Longitud de curva} \index{Longitud de curva}
Para la función  vectorial $\ve{r}(t)$ usaremos la notación
$\overrightarrow{r}(t)$ o ${\bf r}(t)$, indistintamente para referirnos a
este tipo de funciones.

De acuerdo con \cite{marsden_calculo_2013}, la longitud de arco $L$ de la curva $C$ en $\mathbb{R}^3$ determinada por la función vectorial
$\overrightarrow{r}(t)=\left\langle x(t) ,y(t),z(t) \right\rangle$ para $t\in \lbrack a,b]$, está dada por
\begin{equation} \label{curvas}
L=\int_{a}^{b} \left\| \overrightarrow{r}^{\prime }(t)\right\| \, dt.
\end{equation}
La expresión anterior es un caso particular de la fórmula (\ref{intdelinea}), considerando $f(x,y,z)=1$.
Al término $\left\| \overrightarrow{r}^{\prime }(t)\right\|dt$ lo denominaremos elemento diferencial de longitud.
\index{Elemento diferencial!de longitud}
\\
Recordemos que la integral del campo escalar $f\left( x,y,z\right) $ a lo largo de la curva $C$,
está determinada por
$$\int_{C}fds=\int_{a}^{b}f\left( x\left( t\right) ,y\left( t\right) ,z\left(
t\right) \right) \left\| \overrightarrow{r}^{\prime }\left( t\right)\right\| dt.$$

Presentaremos enseguida algunos ejemplos de superficies, su curva de intersección y sus longitud.

\subsection{Un paraboloide y un plano}
La ecuación $-x+y+z=1$ representa un plano y la expresión $2z=x^{2}+y^{2}$
representa un paraboloide. \label{curvas1}

\begin{figure}[H]
\begin{center}
  \caption{Gráfico de $2z=x^{2}+y^{2}$ y $-x+y+z=1$}
  \includegraphics[height=6cm]{Img/72a.pdf} \hskip0.5cm
  \includegraphics[height=6cm]{Img/72b.pdf}
  \fuentepropia
\end{center}
\end{figure}

Con las siguientes instrucciones, se obtienen las figuras anteriores:

\begin{tcolorbox}[title=\bf Instrucción \verb"Mathematica"]
\begin{Verbatim}[formatcom=\letraverbatim]
a=ContourPlot3D[{2*z==x^2+y^2,-x+y+z==1},{x,-5,4},
{y,-4,5},{z,0,6},ContourStyle->{Orange,Cyan},
PlotPoints->60,Mesh->False,Ticks->{{-3,0,3},
{-3,0,3},{1,2,4,5}}];
b=ParametricPlot3D[{1+2*Cos[t],-1+2*Sin[t],3+2*Cos[t]
-2*Sin[t]},{t,0,2*Pi},
PlotStyle->{Black,Thickness->0.016}];
Show[a,b]
  \end{Verbatim}
\end{tcolorbox}

Al remplazar $z$ de la ecuación del paraboloide en la ecuación del plano se obtiene: $\left(\frac{x^{2}+y^{2}}{2}\right)-x+y=1$, que se reescribe como
$\left( x-1\right) ^{2}+\left( y+1\right) ^{2}=4$. Esta última expresión representa una circunferencia con centro en el punto $\left( 1,-1\right) $ y
radio $2$ en el plano $xy$. Esta curva se puede parametrizar en la forma
$x-1=2\cos (t) $, $y+1=2\sin (t) $.
Al sustituir $z$ en la ecuación del paraboloide, se obtiene:
\[
z = \frac{(1+2\cos(t))^{2}+\left(-1+2\sin(t)\right) ^{2}}{2}= 3+2\cos(t) -2\sin (t) .
\]
Por tanto, la función vectorial
$$\overrightarrow{r}(t)=\left\langle 1+2\cos (t) ,-1+2\sin \left(
t\right) ,3+2\cos (t) -2\sin (t) \right\rangle ,\,\,
t\in \lbrack 0,2\pi ]$$
describe la curva de intersección entre el paraboloide y el plano.
Además $\left\| \overrightarrow{r}^{\prime }(t)\right\| = 2\sqrt{2\sin (t) \cos (t) +2}$,
la longitud de la curva está dada por
$$\int_{0}^{2\pi }2\sqrt{2\sin (t) \cos (t) +2}dt\approx 17.475.$$
El cálculo anterior se puede evaluar numéricamente como sigue:

\begin{tcolorbox}[title=\bf Instrucción \verb"Mathematica"]
\begin{Verbatim}[formatcom=\letraverbatim]
r={1+2*Cos[t],-1+2*Sin[t],3+2*cos[t]-2Sin[t]}
A=Assuming[t>=0,Simplify[Norm[D[r,t]]]]
NIntegrate[A,{t,0,2*Pi}]
\end{Verbatim}
\end{tcolorbox}

\subsection{Un paraboloide y un cilindro}
Las ecuaciones $x-y^{2}+4y-z^{2}+2z=2 $ y $y^{2}-4y+4+z^{2}-4z=0$, se reescriben como
$x+3=(y-2)^2+(z-1)^2$ y $(z-2)^2+(y-2)^2=4$.

Estas expresiones representan un paraboloide
con vértice en el punto $(-3,2,1)$, que abre hacia la parte positiva del eje $x$ y un cilindro de
radio $2$ cuyo eje de simetría es una recta paralela al eje $x$. \label{curvas2}

\begin{figure}[H]
\centering
\caption{Gráfico de $x-y^{2}+4y-z^{2}+2z=2$ y $y^{2}-4y+4+z^{2}=4z$}
\includegraphics[width=0.4\textwidth]{Img/70a.pdf} \hspace{5mm}
\includegraphics[width=0.4\textwidth]{Img/70b.pdf}
\fuentepropia
\end{figure}

Las instrucciones para conseguir el gráfico son las siguientes:

\begin{tcolorbox}[title=\bf Instrucción \verb"Mathematica"]
\begin{Verbatim}[formatcom=\letraverbatim]
a=ContourPlot3D[{x+3==(y-2)^2+(z-1)^2,(z-2)^2+(y-2)^2
==4},{x,-3,7},{y,-2,6},{z,-2.2,4.5},Mesh->False,
ContourStyle->{Yellow,RGBColor[0.8,0.3,0.2]},
PlotPoints->60];
b=ParametricPlot3D[{2+4*Sin[t],2+2*Cos[t],
2+2*Sin[t]},
{t,0,2*Pi},PlotStyle->{Black,Thickness->0.014},
PlotPoints->80];
Show[a,b]
\end{Verbatim}
\end{tcolorbox}

Eliminando de las ecuaciones a $y$, se obtiene que $x=2z-2$, que corresponde a la ecuación
de un plano, la curva de intersección está contenida en este. A partir de la ecuación del cilindro, se puede establecer que $y=2+2\cos(t)$ y $z=2+2\sin(t)$
con lo cual $x=2(2+2\sin(t))-2=2+4\sin(t)$.
Entonces, la función vectorial
\[
{\bf r}(t) = \langle 2+4\sin(t), 2+2\cos(t),2+2\sin(t)\rangle,
\]
con $t \in [0,2\pi], $ parametriza la curva de intersección del paraboloide y el cilindro.
A partir de esta función, se sigue que la magnitud del vector tangente es
\[
\left\| {\bf r}^{\prime }(t)\right\| = 2\sqrt{4\cos ^{2}(t) +1}= 2\sqrt{3+2\cos(2t);}
\]
entonces, la longitud de la curva se determina por 
\[
\int_{0}^{2\pi }2\sqrt{3+2\cos(2t)}\,dt \approx 21.081.
\]
Este cálculo se puede obtener numéricamente con las siguientes sentencias:

\begin{tcolorbox}[title=\bf Instrucción \verb"Mathematica"]
\begin{Verbatim}[formatcom=\letraverbatim]
r={2+4*Cos[t],2+2*Cos[t],2+2*Sin[t]};
A=Assuming[t>=0,Simplify[Norm[D[r,t]]]]
NIntegrate[A,{t,0,2*Pi}]
\end{Verbatim}
\end{tcolorbox}

\subsection{Un cilindro y un paraboloide hiperbólico}
El cilindro con ecuación $x^{2}+y^{2}=4$ y el paraboloide hiperbólico con ecuación
$z=x^{2}-y^{2}$ se intersecan en una curva como lo muestra las siguientes figuras.

\begin{figure}[H]
\centering
\caption{Gráfico de $x^{2}+y^{2}=4$ y $z=x^{2}-y^{2}$}
\label{fig:67}
\includegraphics[width=0.4\textwidth]{Img/67a.pdf} \hspace{5mm}
\includegraphics[width=0.4\textwidth]{Fig/301.pdf}
\fuentepropia
\end{figure}

A partir de la ecuación del cilindro se puede definir $x=2\cos \left(t\right) $ y
$\,y=2\sin (t) $ para reemplazar en la ecuación del
paraboloide hiperbólico
$$z=\left( 2\cos (t) \right)^{2}-\left(2\sin(t)\right) ^{2}= 4\cos(2t) ,$$
es decir, la función vectorial
$\overrightarrow{r}\left(t\right) =\left\langle 2\cos (t) ,2\sin (t) ,4\cos\left( 2t\right) \right\rangle $,
para $t\in \left[ 0,2\pi \right] $, define
la curva de intersección de las superficies se\~{n}aladas.
\\
A partir de $\overrightarrow{r}(t) $, se obtiene
$\overrightarrow{r}^{\prime }(t) =\left\langle -2\sin \left(
t\right) ,2\cos (t) ,-8\sin \left( 2t\right) \right\rangle $  y
{\small $$\left\| \overrightarrow{r}^{\prime }(t) \right\| =
2\sqrt{64\cos ^{2}(t) +1-64\cos ^{4}(t)}= 2\sqrt{9-8\cos \left( 4t\right)}.$$
}
Con lo cual, la longitud de esta curva está dada por
{\small $$\int_{0}^{2\pi }2\sqrt{9-8\cos \left( 4t\right) }dt \approx 35.\, 2571.$$
}
El gráfico y la evaluación numérica de esta integral se obtienen en \linebreak
\verb"Mathematica" mediante el uso se las siguientes instrucciones:

\begin{tcolorbox}[title=\bf Instrucción \verb"Mathematica"]
\begin{Verbatim}[formatcom=\letraverbatim]
r={2*Cos[t],2*Sin[t],4*Cos[2*t];Show[ContourPlot3D[
{x^2+y^2==4,z==x^2-y^2},
{x,-4,4},{y,-4,4},{z,-5,5},
Mesh->False,PlotPoints->50,
ContourStyle->{Orange,RGBColor[0.1,0.6,0.2]}],
ParametricPlot3D[{2*Cos[t],2*Sin[t],4*Cos[2*t]},
{t,0,2*Pi},PlotStyle->{Black,Thickness[0.016]}]]
A=Assuming[t>=0,Simplify[Norm[D[r,t]]]]
NIntegrate[A,{t,0,2*Pi}]
\end{Verbatim}
\end{tcolorbox}
%\vspace*{-\baselineskip}
%

\subsection{Un cilindro parabólico y un paraboloide}
Las ecuaciones $z=x^2+y^2$ y $z=5-y^2$, que representan un paraboloide y un cilindro parabólico, respectivamente.

\begin{figure}[H]
\centering
\caption{Gráfico de $z=x^2+y^2$ y $z=5-y^2$}
\includegraphics[width=0.4\textwidth]{Img/50a.pdf} \hspace{5mm}
\includegraphics[width=0.4\textwidth]{Img/50b.pdf}
\fuentepropia
\end{figure}

Las figuras de estas superficies y su curva de intersección se consiguen con las siguientes instrucciones:

\begin{tcolorbox}[title=\bf Instrucción \verb"Mathematica"]
\begin{Verbatim}[formatcom=\letraverbatim]
A=ContourPlot3D[{z==5-y^2,z==x^2+y^2},{x,-3,3},
{y,-3,3},{z,0,5},
PlotPoints->90,ContourStyle->{Cyan,Orange},
Mesh->False]r={Sqrt[5]*Cos[t],Sqrt[5/2]*Sin[t],
(5/2)*(1+Cos[t]^2)};
B=ParametricPlot3D[r,{t,0,2Pi},
PlotStyle->{Black,
Thickness->0.012}];
Show[A,B]
\end{Verbatim}
\end{tcolorbox}

Con las ecuaciones $z=x^{2}+y^{2}$ y $z=5-y^{2}$, se pueden restar y se obtiene
$x^{2}+2y^{2}=5;$ esta expresión representa una elipse con centro en el origen y alargada en el sentido de las $x$.
La curva de intersección está dada en el plano $xy$ por
$x=\sqrt{5}\cos t$, $y=\sqrt{\frac{5}{2}}\sin t$, donde $t\in \lbrack 0,2\pi ]$.
Para hallar la coordenada $z$ se sustituye en la ecuación del paraboloide
$$z=\left( \sqrt{5}\cos t\right) ^{2}+\left( \sqrt{\frac{5}{2}}
\sin t\right) ^{2}= \frac{5}{2}\cos ^{2}t+\frac{5}{2},$$
o también en la ecuación del cilindro
$z=5-\left( \sqrt{\frac{5}{2}}\sin t\right) ^{2}= \frac{5}{2}\cos^{2}t+\frac{5}{2};$
por lo tanto, la curva de intersección está determinada por
$$\overrightarrow{r}(t)=\left\langle \sqrt{5}\cos t,\sqrt{\frac{5}{2}}\sin t,%
\frac{5}{2}\cos ^{2}t+\frac{5}{2}\right\rangle ,\,\,t\in \lbrack 0,2\pi ].$$
Además
$\left\| \overrightarrow{r}^{\prime}(t)\right\| =\frac{\sqrt{10}}{4}\sqrt{11-2\cos(2t)-5\cos(4t)}$,
entonces la longitud de la curva está dada por
$$L=\int_{0}^{2\pi}\frac{\sqrt{10}}{4}\sqrt{11-2\cos(2t)-5\cos(4t)} \,dt \approx 16.192.$$
La integral anterior se evaluó numéricamente mediante las siguientes tres instrucciones:
%\vspace*{-1.2\baselineskip}
%
\begin{tcolorbox}[title=\bf Instrucción \verb"Mathematica"]
\begin{Verbatim}[formatcom=\letraverbatim]
r={Sqrt[5]*Cos[t],Sqrt[5/2]*Sin[t],(5/2)*(1+Cos[t]^2)};
A=Assuming[t>=0,Simplify[Norm[D[r,t]]]]
NIntegrate[A,{t,0,2Pi}]
\end{Verbatim}
\end{tcolorbox}
%\vspace*{-1.2\baselineskip}
%
\subsection{Dos cilindros parabólicos}

La región del espacio acotado por las figuras de las expresiones
$x=z^{2}$ y $x=4y-y^{2}$ \label{curvas3} que corresponden a dos cilindros parabólicos,
junto con la curva de intersección se muestran a continuación.

\begin{figure}[H]
\centering
\caption{Gráfico de $x=z^{2}$ y $x=4y-y^{2}$}
\includegraphics[width=0.4\textwidth]{Img/53a.pdf} \hspace{5mm}
\includegraphics[width=0.4\textwidth]{Img/53e.pdf}

\includegraphics[width=0.4\textwidth]{Img/53f.pdf}
\fuentepropia
\end{figure}

La figura se obtiene con ayuda de las siguientes instrucciones:

\begin{tcolorbox}[title=\bf Instrucción \verb"Mathematica"]
\begin{Verbatim}[formatcom=\letraverbatim]
A=ContourPlot3D[{x==z^2,x==4*y-y^2},{x,-1,5},{y,-1,5},
{z,-3,3},Mesh->False,PlotPoints->50,ContourStyle->
{Blend[{Green,Yellow},2/3],RGBColor[1,0.3,0.3]}];
B=ParametricPlot3D[{4*Sin[t]^2,2+2*Cos[t],2*Sin[t]},
{t,0,2*Pi},PlotStyle->{Black,Thickness[0.016]}];
Show[A,B]
\end{Verbatim}
\end{tcolorbox}

Con estas expresiones se puede reemplazar el valor de $x$ para determinar su proyección en el plano $yz$.
De este procedimiento se obtiene la ecuación  $(y-2)^2+z^2=4$, que corresponde a una circunferencia de radio $2$ centrada en el punto $(2,0)$.  En este caso puede hacerse $y=2+2\cos(t)$ y $z=2\sin(t)$.
La expresión para $x$ se obtiene sustituyendo $y$ y $z$ en cualquiera de las expresiones que determinan los cilindros.
Es decir, $x=4(2+2\cos t) -( 2+2\cos t) ^{2}= 4-4\cos^{2}(t) =4\sin ^{2}(t) $ o también
$x=( 2\sin t) ^{2}= 4\sin ^{2}(t) $.
Por lo tanto, la curva de intersección está dada por la siguiente función vectorial
$${\bf r}(t)=\langle4\sin^{2}(t),2+2\cos(t),2\sin(t)\rangle, \, t \in [0,2\pi].$$
La longitud de esta curva se calcula integrando la norma de la derivada de la función vectorial, es decir
$${\bf r}'(t)= \langle 8\sin(t)\cos(t),-2\sin (t),2\cos(t)\rangle .$$
Su norma está dada por
$$\|{\bf r}^{\prime}(t)\| = 2\sqrt{16\cos^{2}(t)+1-16\cos^{4}(t)}= 2\sqrt{3-2\cos(4t)},$$
este término se integra numéricamente
$$ \int_{0}^{2\pi}2\sqrt{3-2\cos(4t)}dt \approx 21.081.$$
La evaluación se consigue con ayuda de las siguientes instrucciones:

\begin{tcolorbox}[title=\bf Instrucción \verb"Mathematica"]
\begin{Verbatim}[formatcom=\letraverbatim]
r={4*Sin[t]}^2,2+2*Cos[t],2*Sin[t;]
A=Assuming[t>=0,Simplify[Norm[D[r,t]]]]
NIntegrate[A,{t,0,2Pi}]
\end{Verbatim}
\end{tcolorbox}

\subsection{Dos paraboloides}
Las ecuaciones $8y=9-5x^{2}-6z^{2}$ y $16y=4z^{2}-x^{2}$
representan gráficamente un paraboloide y un paraboloide hiperbólico, respectivamente.

\begin{figure}[H]
\centering
\caption{Gráfico de $8y=9-5x^{2}-6z^{2}$ y $16y=4z^{2}-x^{2}$}
\includegraphics[width=0.4\textwidth]{Img/71c.pdf} \hspace{5mm}
\includegraphics[width=0.4\textwidth]{Img/71a.pdf}
\fuentepropia
\end{figure}

Los gráficos se consiguen ejecutando las siguientes instrucciones:

\begin{tcolorbox}[title=\bf Instrucción \verb"Mathematica"]
\begin{Verbatim}[formatcom=\letraverbatim]
a=ContourPlot3D[{8*y+5*x^2+6*z^2==9,16*y==4*z^2-x^2},
{x,-3,3},{y,-1,2},{z,-2,2},Mesh->False,PlotPoints->60,
ContourStyle->{Yellow,RGBColor[0.6,0.2,0.2]},
Ticks->{{-2,0,2},{-1,0,1},{-2,-1,1,2}}];
b=ParametricPlot3D[{Sqrt[2]*Cos[t],(5-13*Cos[2*t])/64,
3*Sqrt[2]/4*Sin[t]},{t,0,2*Pi},PlotPoints->80,
PlotStyle->{Black,Thickness->0.014}];
Show[a,b]
\end{Verbatim}
\end{tcolorbox}

Al eliminar $y$ de estas expresiones se obtiene la ecuación
$9x^{2}+16z^{2}=18$, que corresponde a una elipse en el plano $xz$, esta curva se puede
parametrizar en la forma $x=\sqrt{2}\cos (t) $, $z=\frac{3\sqrt{2}}{4}\sin \left(
t\right) $.
Al reemplazar estas expresiones en cualquiera de las ecuaciones
del paraboloide o del paraboloide hiperbólico se obtiene:
$$8y=9-5x^{2}-6z^{2}= 9-10\cos ^{2}(t) -\frac{27}{4}\sin ^{2}(t)= \frac{9}{4}-\frac{13}{4}\cos^{2}( t),$$
entonces, $y=\frac{9}{32}-\frac{13}{32}\cos ^{2}(t) $.
Por lo tanto, la curva de intersección se parametriza como
$$\overrightarrow{r}(t) =\left\langle \sqrt{2}\cos \left(
t\right) ,\frac{9}{32}-\frac{13}{32}\cos ^{2}(t) ,\frac{3\sqrt{2}}{4}\sin (t) \right\rangle ,\,\, t\in \lbrack 0,2\pi ].$$
Es sencillo ver que
$\left\| \overrightarrow{r}^{\prime }(t)\right\|= \frac{1}{16}\sqrt{512-55\cos ^{2}(t) -169\cos ^{4}(t)} $,
con lo cual, la longitud de la curva es
{\small $$\int_{0}^{2\pi }\frac{1}{16}\sqrt{512-55\cos ^{2}(t) -169\cos^{4}(t) }dt \approx 8.019.$$}

Las siguientes instrucciones permiten hacer este cálculo numéricamente:

\begin{tcolorbox}[title=\bf Instrucción Mathematica]
\begin{Verbatim}[formatcom=\letraverbatim]
r={Sqrt[2]*Cos[t],9/32-13/32*Cos[t]^2,3*Sqrt[2]/
4*Sin[t]};
A=Assuming[t>=0,Simplify[Norm[D[[r,t]]]]
NIntegrate[A,{t,0,2*Pi}]
\end{Verbatim}
\end{tcolorbox}

\subsection{Dos cilindros elípticos y un paraboloide}

Los cilindros elipticos con
ecuaciones $4x^2+y^2=36$ y $x^2+4y^2=4$, poseen al eje $z$ como eje de simetría.
Estos se pueden proyectar en el plano $xy$, como dos elipses que se parametrizan respectivamente como
$\overrightarrow{u}(t)=\left\langle 3\cos (t) ,6\sin (t) \right\rangle $
y  $\overrightarrow{v}(t)=\left\langle 2\cos (t) ,\sin (t) \right\rangle $,
con  $t\in\lbrack 0,2\pi ]$.

Si estos cilindros los intersecan con el paraboloide determinado por la ecuación
$z=6-\left(x-1\right)^2-y^2$,
se obtienen dos curvas, cuya coordenada $z$ está dada, respectivamente por
{\small
\begin{eqnarray*}
z&=&6-\left( 3\cos (t) -1\right) ^{2}-\left( 6\sin (t) \right) ^{2}= -31+27\cos ^{2}(t) +6\cos (t) \\
z &=& 6-\left( 2\cos (t) -1\right) ^{2}-\left( \sin (t) \right) ^{2}= 4-3\cos ^{2}(t) +4\cos (t).
\end{eqnarray*}
}

\begin{figure}[H]
%
\centering
\caption{Gráfico de $4x^2+y^2=36$ y $x^2+4y^2=4$}
\begin{minipage}[t]{0.34\linewidth}
\centering
\includegraphics[width=0.69\linewidth]{Img/1f.pdf}
\end{minipage}%
\hspace{0.2em}
\begin{minipage}[t]{0.34\linewidth}
\centering
\includegraphics[width=0.69\linewidth]{Fig/78.pdf}
\end{minipage}
\fuentepropia
\end{figure}

Con las siguientes instrucciones se consigue graficar estas superficies:

\begin{tcolorbox}[title=\bf Instrucción \verb"Mathematica"]
\begin{Verbatim}[formatcom=\letraverbatim]
ContourPlot3D[{z==6-(x-1)^2-y^2,4*x^2+y^2==36,
x^2+4*y^2==4},{x,-6,10},{y,-7,7},{z,-35,6},
ContourStyle->{Red,Yellow,Green},PlotPoints->80]
\end{Verbatim}
\end{tcolorbox}

En el espacio estas las dos curvas de intersección se pueden parametrizar como
\begin{eqnarray*}
\overrightarrow{\alpha }(t) &=& \left\langle 3\cos (t) ,6\sin
(t) ,-31+27\cos ^{2}(t) +6\cos (t) \right\rangle \\
\overrightarrow{\beta }(t) &=& \left\langle 2\cos (t) ,\sin \left(
t\right) ,4-3\cos ^{2}(t) +4\cos (t) \right\rangle,
\end{eqnarray*}
con $t \in [0,2\pi]$.
La longitud de dichas curvas se puede determinar con el
siguiente procedimiento:
\begin{eqnarray*}
\overrightarrow{\alpha }^{\prime }(t)&=& \left\langle -3\sin (t),6\cos (t) ,-54\cos (t) \sin (t)
-6\sin (t) \right\rangle \\
\left\| \overrightarrow{\alpha }^{\prime }(t)\right\|
&=&  3\sqrt{5+\cos (t) \left( 323\cos(t) -324\cos ^{3}(t) +72\sin ^{2}(t)\right),}
\end{eqnarray*}
con lo cual, la longitud de la curva mayor está dada por
$$\int_{0}^{2\pi}3\sqrt{5+\cos (t) \left( 323\cos(t) -324\cos ^{3}(t) +72\sin ^{2}(t)\right) }  \,  dt\approx 115. 268;$$
por otro lado,
\begin{eqnarray*}
\overrightarrow{\beta }^{\prime }(t)&=& \left\langle -2\sin \left(
t\right) ,\cos (t) ,6\cos (t) \sin (t)-4\sin (t) \right\rangle \\
\left\| \overrightarrow{\beta }^{\prime }(t)\right\|
&=& \sqrt{20+\cos (t) \left( 17\cos (t)-36\cos ^{3}(t) -48\sin ^{2}(t) \right).}
\end{eqnarray*}
Entonces, la longitud de la segunda curva está dada por
$$\int_{0}^{2\pi}\sqrt{20+\cos (t) \left( 17\cos (t)-36\cos ^{3}(t) -48\sin ^{2}(t) \right)} \, dt\approx 21.161.$$
En los dos casos, estas integrales se evaluaron en \verb"Mathematica"
usando las siguientes instrucciones:

\begin{tcolorbox}[title=\bf Instrucción \verb"Mathematica"]
\begin{Verbatim}[formatcom=\letraverbatim]
r={3*Cos[t],6*Sin[t]};
z=6-(r[[1]]-1)^2-r[[2]]^2;
r={r[[1]],r[[2]],z};
NIntegrate[Norm[D[r,t]],{t,0,2*Pi}]
\end{Verbatim}
\end{tcolorbox}

\begin{tcolorbox}[title=\bf Instrucción \verb"Mathematica"]
\begin{Verbatim}[formatcom=\letraverbatim]
r={2*Cos[t],Sin[t]};
z=6-(r[[1]]-1)^2-r[[2]]^2;
r={r[[1]],r[[2]],z};
NIntegrate[Norm[D[r,t]],{t,0,2*Pi}]
\end{Verbatim}
\end{tcolorbox}

\subsection{Un elipsoide y un hiperboloide} \index{Hiperboloide}

Las ecuaciones $2x^2+y^2=z^2+12$ y $x^2+y^2+8z^2=30$ representan un hiperbólico de un manto y un elipsoide, respectivamente.

\begin{figure}[H]
\begin{center}
\caption{Gráfico de $2x^2+y^2=z^2+12$ y $x^2+y^2+8z^2=30$}
\includegraphics[height=4.5cm]{Img/87a.pdf}
\includegraphics[height=4.5cm]{Img/87b.pdf}
\fuentepropia
\end{center}
\end{figure}

La figura anterior se obtiene con las siguientes instrucciones:

\begin{tcolorbox}[title=\bf Instrucción \verb"Mathematica"]
\begin{Verbatim}[formatcom=\letraverbatim]
A=ContourPlot3D[{z^2+12==2*x^2+y^2,x^2+y^2+8*z^2==30},
{x,-6,6},{y,-6,6},{z,-3,3},Mesh->False,PlotPoints->40,
ContourStyle->{Directive[Orange,Opacity[0.8]],
Directive[Cyan,Opacity[0.76]]},BoxRatios->{6,5,4}];
B=ParametricPlot3D[{{3*Sqrt[2]*Sinh[t],Sqrt[2]*Sqrt[6
+Cosh[t]^2-18*Sinh[t]^2],Sqrt[2]*Cosh[t]},
{3*Sqrt[2]*Sinh[t],-Sqrt[2]*Sqrt[6+Cosh[t]^2-
18*Sinh[t]^2],Sqrt[2]*Cosh[t]}},
{t,-ArcCosh[Sqrt[27/17]],ArcCosh[Sqrt[27/17]]},
PlotStyle->{Directive[Black,Thickness->0.015],
Directive[Black,Thickness->0.015]}];
Show[A,B]
\end{Verbatim}
\end{tcolorbox}

Al restar las ecuaciones del elipsoide y del hiperboloide se elimina a $y$,
esto se reduce a $x^{2}-9z^{2}=-18$, que
equivale a $\frac{z^{2}}{2}-\frac{x^{2}}{18}=1;$ esta ecuación representa
una hipérbola y se parametriza como
$z=\sqrt{2}\cosh \left( t\right) $, $x=3\sqrt{2}\sinh \left( t\right) $.
La coordenada $y$ se obtiene reemplazando $x$ y $z$ en cualquiera de las
ecuaciones iniciales (la del elipsoide o la del hiperboloide). Lo cual implica que
$y=\pm \sqrt{48-34\cosh ^{2}t}$.
Para determinar el dominio, puede verse que si $48-34\cosh ^{2}t=0$, entonces la solución está dada por
$a={\rm arccosh}\left( \frac{2\sqrt{102}}{17}\right) \approx 0.7066$.

Por lo descrito anteriormente, la curva de intersección se parametriza por
$$\overrightarrow{r}\left( t\right) =\left\langle 3\sqrt{2}\sinh \left(
t\right) ,\sqrt{48-34\cosh ^{2}t},\sqrt{2}\cosh \left( t\right)
\right\rangle ,$$
con $t \in[-a,a]$, con lo cual, la longitud de la curva está dada por
$$\bigint_{-a}^{a}\sqrt{\frac{2\left(32\cosh ^{2}t-119\cosh^{4}t+24\right)}{17\cosh ^{2}t-24}}dt\approx 30.5672.$$
Los cálculos que permiten obtener la longitud de la curva
se obtienen ejecutando en \verb"Mathematica" las siguientes instrucciones:

\begin{tcolorbox}[title=\bf Instrucción \verb"Mathematica"]
\begin{Verbatim}[formatcom=\letraverbatim]
z=Sqrt[2]*Cosh[t],Sqrt[2]*Cosh[t]}
x=3*Sqrt[2]*Sinh[t];
y=Sqrt[2]*Sqrt[6+Cosh[t]^2-
18*Sinh[t]^2];
a=ArcCosh[Sqrt[27/17]];
A=Assuming[t>=0,Simplify[Norm[D[{r,t},]]]];
NIntegrate[2*A,{t,-a,a}]
\end{Verbatim}
\end{tcolorbox}

\subsection{Una esfera y dos cilindros perpendiculares}

La esfera centrada en el origen de radio $2$ se interseca con los cilindros cuyas ecuaciones son $x^{2}+y^{2}=2x$ y $x^{2}+z^{2}=2x$.
Hallaremos la curva de intersección y de esta calcularemos su longitud de arco. La siguiente figura muestra dicha situación.

\begin{figure}[H]
\centering
\caption{Gráfico de $x^{2}+y^{2}=2x$ y $x^{2}+z^{2}=2x$}
\begin{minipage}[t]{0.38\linewidth}
\centering
\includegraphics[width=0.75\linewidth]{Img/6a.pdf}
\end{minipage}%
\hspace{0.1em}
\begin{minipage}[t]{0.40\linewidth}
\centering
\includegraphics[width=0.80\linewidth]{Fig/79.pdf}
\end{minipage}
\fuentepropia
\end{figure}

Ejecutando las siguientes instrucciones se obtiene la figura:

\begin{tcolorbox}[title=\bf Instrucción \verb"Mathematica"]
\begin{Verbatim}[formatcom=\letraverbatim]
ContourPlot3D[{x^2+y^2+z^2==4,x^2+y^2==2*x,
x^2+z^2==2*x},{x,-2,2},{y,-2,2},{z,-2,2},
Mesh->False,ContourStyle->{Orange,Yellow,Gray}]
\end{Verbatim}
\end{tcolorbox}

El cilindro $x^{2}+y^{2}=2x$, se reescribe como $\left( x-1\right)^{2}+y^{2}=1$,
la curva de intersección se parametriza al hacer $x=1+\cos (t)$, $y=\sin (t)$, y para hallar $z$
es necesario reemplazar en la ecuación de la esfera, de donde se obtiene que

\[x^{2}+y^{2}+z^{2}= \left( 1+\cos (t) \right)
^{2}+\sin ^{2}(t) +z^{2}= 4,\]
es decir, $2+2\cos(t) +z^{2}=4$,
de donde $z=\pm \sqrt{2-2\cos (t)}$,
entonces la función

$${\bf r}_1(t) =\left\langle 1+\cos (t) ,\sin (t),\sqrt{2-2\cos (t) }\right\rangle ,\,\, t\in \lbrack 0,2\pi ],$$

parametriza la curva de intersección. Es de anotar que son dos porciones que se obtienen por la raiz positiva y la raiz negativa.

Para la expresión  $x^{2}+z^{2}=2x$ (que representa el otro cilindro), el tratamiento es similar, obtiendose la función
$${\bf r}_2(t) =\left\langle 1+\cos (t),\sqrt{2-2\cos (t)},\sin (t)\right\rangle ,$$
donde $t\in \lbrack 0,2\pi ]$.
Por la simetría de la curva y las ecuaciones que la definen
consideraremos solamente ${\bf r}_1(t)$, entonces la longitud ($L$) de la curva es cuatro veces la integral del
término
$$\left\Vert {\bf r}_1^{\prime }(t) \right\Vert =\left\Vert \left\langle -\sin t,\cos t,\frac{1}{2}\frac{\sqrt{2}\sin t}{\sqrt{1-\cos t}}\right\rangle \right\Vert=\frac{1}{2}\sqrt{2}\sqrt{\cos t+3} ,$$
de donde
$$L=4\int_{0}^{2\pi }\frac{1}{2}\sqrt{2}\sqrt{\cos t+3}\,dt \approx 30.562.$$

Este valor se obtuvo en \verb"Mathematica" utilizando las siguientes instrucciones:

\begin{tcolorbox}[title=\bf Instrucción \verb"Mathematica"]\begin{Verbatim}[formatcom=\letraverbatim]
r={1+Cos[t],Sin[t],Sqrt[2-2Cos[t]]};
A=Assuming[t>=0,Simplify[Norm[D[r,t]]]];
NIntegrate[4*A,{t,0,2Pi}]
\end{Verbatim}
\end{tcolorbox}

\subsection{Un cilindro cardióidico y un paraboloide} \label{cilcardiopar}

Consideremos el parabóloide con ecuación $3z=9-(x+2)^{2}-(y+1)^{2}$ y el cilindro cardióidico
$r=1-\cos (t)$. En dos dimensiones, el cardioide se  parametriza como
$\overrightarrow{r}(t)=\left\langle (1-\cos (t))\cos (t),(1-\cos (t))\sin(t)\right\rangle$,
donde $t\in [0,2\pi]$.

\begin{figure}[H]
\centering
\caption{$r=1-\cos t$, $3z=4-4x-x^2-2y-y^2$}
\includegraphics[width=0.4\textwidth]{Img/58a.pdf} \hspace{5mm}
\includegraphics[width=0.4\textwidth]{Img/58b.pdf}
\fuentepropia
\end{figure}

Sobre el paraboloide se cumple que
$z=\frac{9}{2}-2 \sin (t)+\sin (2 t)-2 \cos (t)+\frac{3}{2} \cos (2 t)$.
Por lo tanto, para $t\in [0,2\pi]$, la curva de intersección del paraboloide y el cilindro cardioidico se parametriza por medio de la sigui\-ente función:
\begin{multline*}
\overrightarrow{r}(t)=\big\langle(1-\cos t)\cos t,(1-\cos t)\sin t,\frac{9}{2}-2\sin t+\sin(2t) \\
-2\cos t+\frac{3}{2}\cos(2t)\big\rangle.
\end{multline*}

Puede verse también, que:
\begin{multline*}
\Vert \overrightarrow{r}^\prime(t)\Vert= \frac{\sqrt{2}}{3}\left|\sin\left(\frac{t}{2}\right)\right|\bigg(4(\sin(t)+\sin(2t)+3\sin(3t))+7\cos(t)\\ +6\cos(2t)+5\cos(3t)+32\bigg )^{1/2} ,
\end{multline*}

por tanto:
$$\int_0^{2\pi}\left\Vert \overrightarrow{r}^\prime(t)\right\Vert dt \approx 10.059.$$

Las figuras y los cálculos para determinar la longitud de la curva se obtienen ejecutando las siguientes instrucciones:

\begin{tcolorbox}[title=\bf Instrucción \verb"Mathematica"]\begin{Verbatim}[formatcom=\letraverbatim]
r={(1-Cos[t])*Cos[t],(1-Cos[t])*Sin[t],
1+Cos[t]/3*(3*Cos[t]-2)+2*Sin[t]/3*(Cos[t]-1)};
A=Assuming[t>=0,Simplify[Norm[D[r,t]]]]
NIntegrate[A,{t,0,2*Pi}]
a=ContourPlot3D[{3*z==9-(x+2)^2-(y+1)^2},{x,-5,1},
{y,-4,2},{z,0,3},ContourStyle->{Orange},Mesh->False,
PlotPoints->80];
b=ParametricPlot3D[{r[[1]],r[[2]],u},{t,0,2*Pi},
{u,0,3},Mesh->False,PlotPoints->80,
PlotStyle->{Opacity[0.7],RGBColor[0.1,0.5,0.3]}];
c=ParametricPlot3D[r,{t,0,2*Pi},PlotPoints->80,
PlotStyle->{Black,Thickness->0.015}];
Show[a,b,c]
\end{Verbatim}
\end{tcolorbox}

\subsection{Superficies obtenidas a partir de su curva de intersección}

Para las siguientes funciones vectoriales se construyen dos superficies cuya intersección corresponde
a la curva determinada por la función $\overrightarrow{r}(t)$. Además se halla la longitud ($L$) de la curva para
un intervalo dado.

%

\begin{funciones}[series=ejercicios]
\item \makebox[\linewidth][c]{$\overrightarrow{r}(t)=\la 3t^{3}+t,3t^{2},3t^{3}-t\ra, $ $t \in [-2,2]$.}
\end{funciones}

De las ecuaciones paramétricas se sigue que
$x=3t^{3}+t$, $y=3t^{2}$, $z=3t^{3}-t$.
Entonces, $x-z=2t;$ por tanto, $\left( x-z\right) ^{2}=4t^{2}=\frac{4}{3}\left(
3t^{2}\right) =\frac{4}{3}y$. Además $x+z=6t^{3}=\left( 2t\right) \left( 3t^{2}\right) =\left(
x-z\right) y$.

Por tal razón, utilizaremos las expresiones
$x-z=\frac{4}{3}y$ y $x+z=y\left( x-z\right) $, cuyos gráficos se
presentan a continuación y con éstos también se muestra la curva de intersección.

\begin{figure}[H]
\centering
\caption{Superficies con intersección ${\bf r}(t)=\la 3t^{3}+t,3t^{2},3t^{3}-t\ra$}
\includegraphics[width=0.4\textwidth]{Img/94a.pdf} \hspace{5mm}
\includegraphics[width=0.4\textwidth]{Img/94b.pdf}
\fuentepropia
\end{figure}

Las indicaciones para conseguirlos son las siguientes:

\begin{tcolorbox}[title=\bf Instrucción \verb"Mathematica"]\begin{Verbatim}[formatcom=\letraverbatim]
a=ParametricPlot3D[{3*t^3+t,3*t^2,3*t^3-t},{t,-1,1},
PlotStyle->{Black,Thickness->0.012},PlotPoints->90];
b=ContourPlot3D[{(x-z)^2==4*y/3,x+z==y*(x-z)},
{x,-1,1},{y,-0.5,1},{z,-1,1},Mesh->False,
ContourStyle->{Directive[RGBColor[0.3,0.6,0.3],
Opacity[0.95]],Directive[Orange,Opacity[0.8]]}];
Show[b,a]
\end{Verbatim}
\end{tcolorbox}

A partir de la relación
$\left\| \frac{d}{dt}\la 3t^{3}+t,3t^{2},3t^{3}-t\ra\right\| = \sqrt{2}\left( 9t^{2}+1\right) $, es posible establecer que
la longitud de la arco está dada por
$$L=\int_{-2}^{2}\sqrt{2}\left( 9t^{2}+1\right) dt= 52\sqrt{2}.$$

\begin{funciones}[resume=ejercicios]
\item \makebox[\linewidth][c]{$\overrightarrow{r}(t)=\la \sqrt{t^{3}},\frac{9}{16}t,t^{2}\ra $, $t \in [0,2]$.}
\end{funciones}

A partir de las ecuaciones $x=\sqrt{t^{3}}$, $y=\frac{9}{16}t$, $z=t^{2}$, se sigue que $x^2=t^3$ y $t= \frac{16}{9}y;$
entonces, dos posibles relaciones son $9x^2=16yz$ y $z= \left(\frac{16}{9}y\right)^2$.

Las figuras de estas ecuaciones y la curva de intersección se presentan enseguida:

\begin{figure}[H]
\centering
\caption{Superficies con intersección ${\bf r}(t)=\la \sqrt{t^{3}},\frac{9}{16}t,t^{2}\ra$}
\includegraphics[width=0.4\textwidth]{Img/92a.pdf} \hspace{5mm}
\includegraphics[width=0.4\textwidth]{Img/92b.pdf}
\fuentepropia
\end{figure}

Las instrucciones requeridas para obtener las figuras son las siguientes:

\begin{tcolorbox}[title=\bf Instrucción \verb"Mathematica"]\begin{Verbatim}[formatcom=\letraverbatim]
a=ParametricPlot3D[{t^(3/2),9*t/16, t^2},{t,0,1},
PlotStyle->{Black,Thickness->0.014}];
b=ContourPlot3D[{9*x^2==16*y*z,z==(16*y/9)^2},{x,0,1},
{y,0,1},{z,0,2}, PlotPoints->90,Mesh->False,
ContourStyle->{Directive[Orange,Opacity[0.8]],
Directive[RGBColor[0.5,0.8, 0.5],Opacity[0.85]]}];
Show[b,a]
\end{Verbatim}
\end{tcolorbox}

Además, $\left\| \frac{d}{dt}\la \sqrt{t^{3}},\frac{9}{16}t,t^{2}\ra\right\| = \frac{9}{16}+2t$,
entonces la longitud de la curva está dada por
$$L=\int_0^2 \left ( \frac{9}{16}+2t \right )dt = \frac{41}{8}. $$

\begin{funciones}[resume=ejercicios]
\item \makebox[\linewidth][c] {$\overrightarrow{r}(t)=\la 4t^{3}+2t^{2},\frac{4}{3\sqrt{2}}t,4t^{3}-2t^{2}\ra $, $t \in [-1,1]$.}
\end{funciones}

A partir de las expresiones $x=4t^{3}+2t^{2}$, $y=\frac{4t}{3\sqrt{2}}$, $z=4t^{3}-2t^{2}$, se sigue que
$t=\frac{3\sqrt{2}}{4}y$, $x-z=4t^{2}=4\left( \frac{3y}{4}\sqrt{2}\right)
^{2}= \frac{9}{2}y^{2}$.\\
Además $x+z=8t^{3}=8\left( \frac{3y}{4}\sqrt{2}\right) ^{3}= \frac{27}{4}\sqrt{2}y^{3}$.

Por lo tanto, las figuras de las ecuaciones $x-z= \frac{9}{2}y^{2}$ y $x+z=\frac{27}{4}\sqrt{2}y^{3}$
se muestras a continuación junto con la curva de intersección:

\begin{figure}[H]
\begin{center}
\caption{Superficies con intersección ${\bf r}(t)=\la 4t^{3}+2t^{2},\frac{4}{3\sqrt{2}}t,4t^{3}-2t^{2}\ra$}
\includegraphics[height=4.5cm]{Img/93a.pdf} \hskip0.5cm
\includegraphics[height=4.5cm]{Img/93b.pdf}
\fuentepropia
\end{center}
\end{figure}

Estos figuras se obtienen ejecutando las siguientes instrucciones:

\begin{tcolorbox}[title=\bf Instrucción \verb"Mathematica"]\begin{Verbatim}[formatcom=\letraverbatim]
a=ParametricPlot3D[{4*t^3+2*t^2,4*t/(3*Sqrt[2]),
4*t^3-2*t^2},{t,-2,2},PlotPoints->90,
PlotStyle->{Black,Thickness->0.014}];
b=ContourPlot3D[{x+z==27*Sqrt[2]*y^3/4,x-z==9*y^2/2},
{x,-1,1},{y,-1.5,1},{z,-2,1},PlotPoints->90,
Mesh->False,ContourStyle->{Directive[Orange,
Opacity[0.8]],Directive[RGBColor[0.5,0.8,0.5],
Opacity[0.9]]}];
Show[b,a]
\end{Verbatim}
\end{tcolorbox}

Es sencillo ver que
$ \left\| \frac{d}{dt}\la 4t^{3}+2t^{2},\frac{4t}{3\sqrt{2}},4t^{3}-2t^{2}\ra \right\| = \frac{2}{3}\sqrt{2}\left(
18t^{2}+1\right), $
por lo tanto, la longitud de la curva está dada por
$$L=\int_{-1}^{1}\frac{2}{3}\sqrt{2}\left( 18t^{2}+1\right) dt= \frac{28}{3}\sqrt{2}$$

\begin{funciones}[resume=ejercicios]
\item \makebox[\linewidth][c]
{$\overrightarrow{r}(t) =\la 3t-t^{3},3t^2,3t+t^{3} \ra $, $t \in [-1,1]$.}
\end{funciones}

Para esta función es claro que
$x= 3t-t^{3}$, $y=3t^2$, $z=3t+t^{3}$. De modo que
$x+z=6t$ entonces $(x+z)^2=12y$. Por otro lado, $x^2-z^2=-12t^4$, entonces, $x^2-z^2=-4y^2/3$.
Las figuras de las dos superficies junto con la curva de intersección se muestra a continuación:

\begin{figure}[H]
\centering
\caption{Superficies con intersección ${\bf r}(t)=\la 3t-t^{3},3t^2,3t+t^{3} \ra$}
\includegraphics[width=0.4\textwidth]{Img/91a.pdf} \hspace{5mm}
\includegraphics[width=0.4\textwidth]{Img/91b.pdf}
\fuentepropia
\end{figure}

Las instrucciones para conseguir estas figuras son las siguientes:

\begin{tcolorbox}[title=\bf Instrucción \verb"Mathematica"]\begin{Verbatim}[formatcom=\letraverbatim]
a=ParametricPlot3D[{3*t-t^3,3*t^2,3*t+t^3},
{t,-1.15,1.15},PlotStyle->{Black,Thickness->0.017},
PlotPoints->90];
b=ContourPlot3D[{x^2-z^2+4*y^2/3,(x+z)^2==12*y},
{x,-5,5},{y,-5,5},{z,-5,5},
ContourStyle->{Directive[Orange,Opacity[0.7]],
Directive[RGBColor[0.5,0.8,0.5],Opacity[0.6]]}],
Mesh->False;
Show[b,a]
\end{Verbatim}
\end{tcolorbox}

Para hallar $L$ se determina el elemento diferencial de longitud
$$\left\| \frac{d}{dt}\la 3t-t^{3},3t^2,3t+3t^{3} \ra \right\| = 3 \sqrt{2}(t^2+1),$$

por lo tanto, $$L= \int_{-1}^1 3 \sqrt{2}(t^2+1) dt = 8\sqrt{2}.$$

\begin{funciones}[resume=ejercicios]
\item \makebox[\linewidth][c]
{$\overrightarrow{r}(t) =\la \sec t, \ln (\sec(t)+\tan(t)),\tan(t) \ra $, $t \in [-\frac{\pi}{3},\frac{\pi}{3}]$.}
\end{funciones}

A partir de las funciones componentes y de algunas identidades trigono\-métricas conocidas
se pueden establecer relaciones entre ellas;
dos de éstas funciones son: $x^2=1+z^2$ y $y=\ln(x+z)$.
Los figuras de estas dos superficies junto con la curva de intersección se muestra a continuación:

\begin{figure}[H]
\centering
\caption{Superficies con intersección ${\bf r}(t)=\la \sec t, \ln (\sec(t)+\tan(t)),\tan(t) \ra$}
\includegraphics[width=0.4\textwidth]{Img/95a.pdf} \hspace{5mm}
\includegraphics[width=0.4\textwidth]{Img/95b.pdf}
\fuentepropia
\end{figure}

Las instrucciones para conseguir estas figuras son las siguientes:

\begin{tcolorbox}[title=\bf Instrucción \verb"Mathematica"]\begin{Verbatim}[formatcom=\letraverbatim]
a=ParametricPlot3D[{Sec[t],Log[Sec[t]+Tan[t]],Tan[t]},
{t,-Pi/3,Pi/3},PlotStyle->{Thickness[0.017],Black}];
b=ContourPlot3D[{1+z^2==x^2,y==Log[x+z]},{x,-2,3},
{y,-2,2},{z,-2,2},Mesh->False,
ContourStyle->{Directive[Orange,Opacity[0.95]],
Directive[Cyan,Opacity[0.9]]}];
Show[b,a]
\end{Verbatim}
\end{tcolorbox}

Es claro que
\[
\left\| \overrightarrow{r}^{\prime }(t)\right\| =\left\| \frac{d}{dt}\left\langle \sec (t) ,\ln \left( \sec (t) +\tan (t) \right) ,\tan (t) \right\rangle \right\| =\sqrt{2}\sec ^{2}t\text{,}
\]
entonces, la longitud de la curva está dada por
$$L=\int_{-\pi /3}^{\pi /3}\sqrt{2}\sec ^{2}tdt= 2\sqrt{6}.$$
%\end{enumerate}

\begin{ejer}
\normalfont
Para cada una de las siguientes funciones vectoriales, determine al menos dos
superficies cuya intersección esté dada por $\overrightarrow{r}(t)$.
También halle la longitud de la curva para el intervalo dado.
\begin{enumerate}
\item $ \overrightarrow{r}(t)=\la e^{-t}\sin (t),e^{-t},e^{-t}\cos (t) \ra$, $t \in [0,\pi/2]$.
%$$\left\| \frac{d}{dt}\la e^{-t}\sin (t),e^{-t},e^{-t}\cos (t) \ra \right\| = \sqrt{3}e^{-t}$$
\item $ \overrightarrow{r}(t)=\la e^{-t}\cos \left( 2t\right),e^{-t}\sin \left( 2t\right),e^{-t} \ra $,
$t \in [0,\pi/2]$.
%$$\left\| \frac{d}{dt}\la e^{-t}\sin \left( 2t\right),e^{-t},e^{-t}\cos \left( 2t\right) \ra \right\| = \sqrt{6}e^{-t}$$
\item $\overrightarrow{r}(t)=\la e^{t},e^{t}\sin \left(2t\right) ,e^{t}\cos \left( 2t\right) \ra $,
$t \in [0,\pi/2]$
%$$\left\| \frac{d}{dt}\la e^{t},e^{t}\sin \left( 2t\right),e^{t}\cos \left( 2t\right) \ra \right\| = \sqrt{6}e^{t}$$
\item $ \overrightarrow{r}(t)=\la 3t^{4}+t^{2},\frac{4}{3}\sqrt{6}t^{3},3t^{4}-t^{2}\ra $,
$t \in [-1,1]$.
%$$\left\|\frac{d}{dt}\la 3t^{4}+t^{2},\frac{4}{3}\sqrt{6}t^{3},3t^{4}-t^{2}\ra \right\| =  2t\sqrt{2}\left(6t^{2}+1\right) $$
\item $ \overrightarrow{r}(t)=\la t^{2},\frac{6}{\sqrt{3}}t,\ln t^{3}\ra $,
$t \in [1,e]$.
%$$\left\| \frac{d}{dt}\la t^{2},\frac{6}{\sqrt{3}} t,\ln t^{3}\ra \right\|= \frac{2t^{2}+3}{t}$$
\item $\overrightarrow{r}(t) =\la \frac{1}{3}t^{3}-4t+1,-\frac{4}{t}-2t^{2},2\sqrt{2}\left( t+\ln t^{2}\right) \ra $,
$t \in [1,e]$.
\normalfont
\item $\overrightarrow{r}(t) =\la \cosh(2t), 2t, -\sinh(2t) \ra $, $t\in [-1,1]$.
%$$\left\| \frac{d}{dt}\la \frac{1}{3}t^{3}-4t+1,-\frac{4}{t}-2t^{2},\frac{4}{\sqrt{2}}\left( t+\ln t^{2}\right) \ra \right\| = \frac{\left( t^{2}+2\right) ^{2}}{t^{2}}$$
\item $\overrightarrow{r}(t) =\la t \cos (2t),t \sin (2t),7 t \ra $, $t\in [0,2\pi]$.
\item $\overrightarrow{r}(t) =\la t^2\cos (3t),t^2 \sin (3t),t^2 \ra $, $t\in [0,2\pi]$.
\item $\overrightarrow{r}(t) =\la e^t \cosh(t),e^t \sinh(t),e^t \ra $, $t\in [-1,1]$.
\item $\overrightarrow{r}(t) =\la -\frac{1}{t},\frac{2}{3}t^3,4t+5\ra $, $t\in [1,2]$.
\item $\overrightarrow{r}(t) =\la t,\ln (\cot (t)-\csc (t)),\ln(\sin (t)) \ra $, $t\in [\pi/6,5\pi/6]$.

\end{enumerate}
\end{ejer}

\begin{ejer}\textit {Determine una expresión que defina la curva de intersección de las siguientes superficies y halle su longitud.}

\enlargethispage{\baselineskip}

\begin{enumerate}
\item El hiperboloide $x^2+z^2=4y^2+\frac{1}{2}$ y el paraboloide $x^2+z^2=1-y$.

\begin{figure}[H]
\begin{center}
\caption{Gráfico de $x^2+z^2=4y^2+\frac{1}{2}$ y $x^2+z^2=1-y$}
\includegraphics[height=4.5cm]{Img/T_23a.pdf}\hskip0.5cm
\includegraphics[height=4.5cm]{Img/T_23b.pdf}
\fuentepropia
\end{center}
\end{figure}

\item El hiperboloide $x^2+z^2-2x=y^2-2$ y el cilindro $3x^2+2z^2-4z=34$.

\begin{figure}[H]
\begin{center}
\caption{Gráfico de $x^2+z^2-2x=y^2-2$ y $3x^2+2z^2-4z=34$}
\includegraphics[height=4.5cm]{Img/T_24a.pdf}\hskip0.5cm
\includegraphics[height=4.5cm]{Img/T_24b.pdf}
\fuentepropia
\end{center}
\end{figure}

\item La esfera $x^2+y^2+z^2-2x=31$ y el paraboloide $x^2+y^2+2x=4z-1$.

\begin{figure}[H]
\begin{center}
\caption{Gráfico de $x^2+y^2+z^2-2x=31$ y $x^2+y^2+2x=4z-1$}
\includegraphics[height=4.5cm]{Img/T_25a.pdf}\hskip0.5cm
\includegraphics[height=4.5cm]{Img/T_25b.pdf}
\fuentepropia
\end{center}
\end{figure}

\item El cilindro cardioidico con ecuación {\small $r=1+\cos(t)$ y el plano $x+y+3z=1$.} \index{Cilindro cardiodico}

\vspace*{-4mm}
\begin{figure}[H]
\begin{center}
\caption{Gráfico de $r=1+\cos(t)$ y $x+y+3z=1$}
\includegraphics[height=4.5cm]{Img/T_29a.pdf}\hskip0.25cm
\includegraphics[height=4.5cm]{Img/T_29b.pdf}
\fuentepropia
\end{center}
\end{figure}

\item Las esferas
$x^2 + y^2 + z^2 - 2 x = 4$ y $x^2 + y^2 + z^2 - 2z = 4$.

\begin{figure}[H]
\begin{center}
\caption{Gráfico de $x^2 + y^2 + z^2 - 2 x = 4$ y $x^2 + y^2 + z^2 - 2z = 4$}
\includegraphics[height=4.5cm]{Img/T_72a.pdf}\hskip0.5cm
\includegraphics[height=4.5cm]{Img/T_72b.pdf}
\fuentepropia
\end{center}
\end{figure}

\item El cilindro $r=\sin(2\theta)$ y el plano $x+y+3z=1$.

\begin{figure}[H]
\begin{center}
\caption{Gráfico de $r=\sin(2\theta)$ y $x+y+3z=1$}
\includegraphics[height=4.5cm]{Img/97a.pdf}\hskip0.5cm
\includegraphics[height=4.5cm]{Img/97b.pdf}
\fuentepropia
\end{center}
\end{figure}

\item El cilindro $r=\sin(2\theta)$ y el paraboloide $z=1-x^2-y^2$.

\begin{figure}[H]
\begin{center}
\caption{Gráfico de $r=\sin(2\theta)$ y $z=1-x^2-y^2$}
\includegraphics[height=4.5cm]{Img/98a.pdf}\hskip0.5cm
\includegraphics[height=4.5cm]{Img/98b.pdf}
\fuentepropia
\end{center}
\end{figure}
\end{enumerate}
\end{ejer}
